\section{Benchmark scoring and native rewards}\label{benchmark-scoring-and-mechanism-mapping}

The launch rewards submitted \textit{A} artifacts under coordinator-led
scoring. Slop Ninja runs each artifact and two reference models on the
same fresh hidden batch, retains every prediction and failure, and signs the
result package. Every validator checks the published score components,
independently derives the canonical native row under the frozen
\texttt{reward-v1} rules below, and submits it through its own journaled
commit--reveal, in parallel with the other validators (Section~\ref{cohorts}).
Held-out documents, labels, and individual predictions remain inside the
coordinator's evaluation service. Other validators trust the coordinator's
execution and labels while checking the reward calculation themselves. Native
consensus determines the eventual payouts~\cite{btchainconsensus,btemissions}.

\texttt{reward-v1} ports the band mechanism of Bittensor subnet 78
(Umi)~\cite{umisn78}: geometric weights across score bands and within each band,
a baseline that ratchets only on improvement, and a largest-remainder integer
row. Its principal departures are these: artifacts stay private until the
cohort ends; eligibility requires a margin over the reference rather than
subnet 78's 1\% grace below the incumbent; a fallback cutoff, a Slop Ninja rule
with no subnet 78 counterpart, keeps rewards flowing when nobody reaches that
margin; an entrant too close to any published artifact earns nothing, and
near-duplicate entrants share one credit equally; and no record, identity, or
gate carries a chain runtime version. The shared protocol crate implements
these rules as exact computation~\cite{protocolcrate}.

\subsection{Freeze the comparison before execution}

A cohort manifest binds its roster of artifact commitments, the raw-text input
profile and output contract, the reference gallery, the reference models, the
batch design, the human pool, qualification gates, numerical tolerances, the
artifact cap and compute limits (Section~\ref{a-artifact-execution}), retry
policy, the cohort schedule (Section~\ref{cohorts}), and the \texttt{reward-v1}
parameters. The batch design fixes \(n\), the mix, the domain\(\times\)origin
strata, and the sealed batch digest. The human pool enters through its frozen
pool manifest and, for the mainnet stock, its stratum share caps and a
commitment to its secret sampling key (Section~\ref{human-stock}). The manifest also binds the index of
published artifacts and the fingerprint probe versions in use
(Section~\ref{public-copies}). Every model in the comparison, including both
references, receives the same canonical text for each document. Hidden labels
remain outside model inputs.

One model version per miner enters a cohort. On-chain digest commitments and
private uploads have an admission cutoff (Section~\ref{a-artifact-execution});
accepted cohorts have no completion expiry. The credited identity is the
original hotkey and its registration generation, which survives UID changes.
Renaming metadata cannot create new inference content.

The coordinator caches admitted artifacts and the cohort batch before
execution. Runs have no online learning, network access, or undeclared mutable
state, and each document is scored independently of the others in its batch.
Every assigned item receives an outcome record. Infrastructure errors, model
errors, invalid outputs, and unsupported inputs remain distinguishable. The
comparison cannot select the easiest successful subset for each entrant or its
best answer across retries.

\subsection{The cohort batch}\label{cohort-batch}

Each cohort draws a fresh hidden batch, and every entrant and both reference
models score the same documents, so every comparison is paired. The batch has
about 1,000 cases on testnet and 4,000 on mainnet, with 2,000 as the mainnet
floor. Its mix is 20\% human-only, 55\% mixed, and 25\% model-only documents.
Fresh AI-side text, both partial rewrites and model-only mirrors, is generated
for each cohort on serverless services such as Workers AI or hosted model APIs,
under a per-cohort budget, by parallel generation lanes that share the
providers' published rate limits~\cite{cfailimits} (Section~\ref{fan-out}).

Whole source groups are sampled, and the design is self-weighting over
domain\(\times\)origin cells: every document carries equal weight, and no Kish
design loss arises. The mix is therefore part of the estimand. Under 20/55/25,
human-only documents carry about 21\% of the reference's Brier loss, against
43\% under a 40/40/20 mix, so the human share sets the reward's weight on human
false positives. Holding a class-balanced estimand fixed at this human share
would instead raise the paired standard error by 17--21\%. A model's
\emph{headroom} is one minus its score, its remaining Brier loss
(Section~\ref{detection}); margins, band widths, and gains are stated as
fractions of the reference's headroom. Three headroom
figures appear in this paper, each a different estimand computed from the same
\(R_0\) predictions on the consumed Test: 0.0518 is its balanced involvement
Brier under the Test's own pre-registered document weights
(Table~\ref{tab:tier-c}); about 0.053 weights documents equally within each
domain\(\times\)origin cell at the Test's origin mass of about 34/34/32; and
about 0.062 is the self-weighting estimand of a launch batch, under 20/55/25 as
under 40/40/20. Over the eighteen domains of the consumed Test and the testnet
stock, a 1,000-case batch holds six to ten human-only documents per domain,
minimum to median, so per-domain slices serve as diagnostics.

About 200 human-only documents in a 1,000-case batch measure a human
false-positive rate at the frozen 1\% operating point only to about \(\pm1.6\)
points (95\%). The separate pure-human false-positive gate therefore has two
tiers. Every entrant must pass a loose threshold on its own cohort's human-only
documents. A candidate for promotion to reference must also pass a strict
threshold measured on the thousands of human documents in the lookup canary's
fresh Common Crawl draw, which is scored for every such candidate anyway
(Section~\ref{lookup-exposure}). Testnet has neither Common Crawl stock nor
canary anchors, so there the strict tier is measured on the canary's first audit
draw: separate human documents drawn from the same pools as the cohort's
(Section~\ref{human-stock}), at least as many as its human-only documents, and
at that size no more precise than the loose tier. Rates pooled over seven cohorts, about
\(\pm0.6\) points at 1,000 cases and \(\pm0.3\) at 4,000, are reported for
monitoring. Human spans inside mixed documents cannot stand in for these documents:
identical human words score higher inside a mixed document than alone
(Section~\ref{cohort-mix-evidence}), so their in-context false-positive rate is
a separate span measurement (Section~\ref{span-outputs}).

A passage of fresh stock is scored as a pure-human negative at most once. It may also
serve as the source of model rewrites, in its own cohort or later ones, at most
four times and with uses at least seven cohorts apart, with a fresh generator,
seed, voice, and editing operation each time; a use in its own cohort counts
toward the four. A rewrite ledger records each use's generator, voice, and
operation, and the signed manifest binds the use limit, the gap, which counts
cohorts that drew human text, and the set of voices and operations. A mixed document repeats its
source's unreplaced words byte for byte, so a passage is no longer unseen once
used. In steady state only pure-human negatives consume the human stock, about
\(0.2n\) per cohort: 73,000 documents a year at 1,000 cases and 292,000 at
4,000; before then, rewrite sources come from fresh stock or from the same
cohort's negatives. Each document is scored independently, and qualification
requires a document's outputs to be identical whether it is scored alone or
among its batch, so an artifact cannot link a passage to its other appearances,
in the same batch or earlier ones. The dependence that reuse adds between
cohorts is negligible: the within-group correlation is about \(-0.03\).

At 1,000 cases the standard error of a paired score difference between two
independently retrained models is about 3.5\% of the reference's headroom under
20/55/25; it is 2.5\% at 2,000 and 1.7\% at 4,000. These figures were measured
on self-weighting batches drawn from the consumed Test. Separating two models
that differ by 1\% of headroom, the measured training-seed noise floor, would
need about 126,000 documents under the sealed Test's design, where the same
standard error is about 4.0\% of headroom at 1,000 cases, for a two-sided 5\%
test at 80\% power. \texttt{reward-v1} therefore sets its eligibility margin and
band width at about two standard errors; within a band, rank order still sets
the geometric split, so differences below the noise still change pay, and
across a band boundary they can reverse its order (Section~\ref{reward-v1}). Container scoring costs about \$0.03--\$0.34 per
1,000 documents, depending on length and instance size
(Section~\ref{existing-evidence-and-implementation-status}).

\subsection{Human stock}\label{human-stock}

The mainnet human stock is pre-cutoff Common Crawl~\cite{commoncrawl}: pages
from crawls CC-MAIN-2013-20 through CC-MAIN-2022-49 captured before November 30,
2022, the public release of ChatGPT. These pages form a historical-human proxy
stratum, not documented human-only work. The mainnet estimand is therefore
historical human writing against current model output. Pages pass the date
gate, rights exclusions, quality rules, and versioned machine-text screens.
No detector filters individual documents: at least 10\% of each quarterly pool
is scored with \(R_0\) and two earlier detectors, and a domain and capture-era
cell is capped or excluded when the upper 95\% bound of the excess of its later
crawls' flag rate over its earliest crawls' exceeds two points on two of the
three detectors or on the screen counts. Pilot measurements appear in
Section~\ref{cc-pilot-evidence}.

A frozen quarterly pool manifest governs each draw. A secret 256-bit key ranks
WET files, source groups, and the passage taken from each page; the manifest
records the key's SHA-256, and the key is revealed when the pool retires. Random
WET files fill every domain, and draws approximate uniform sampling over unique
documents by weighting each crawl by its new URLs and down-weighting repeated
captures. Versioned domain rules assign each excerpt its cohort domain. A fresh
hand-labelled draw under the current rules found 72\% [64, 79] of documents
usable at the launch shares, with no domain reaching 90\%; the rest are human
text under the wrong domain or register. Domain precision is therefore a stated
limit that sizes the pool, not a launch requirement
(Section~\ref{cc-stock-evidence}). Launch shares of human parents follow
per-domain caps derived from each domain's measured frame: essays 51.25\%, news
26.5\%, forums 7.8\%, reviews 4.16\%, legal text 3.11\%, technical text 2.39\%,
speech 2.3\%, fiction 1.8\%, and email 0.69\%, almost all press releases. Each
capped domain sits under its cap, and essays receive the freed points. A larger
share for business prose needs its own measured frame outside guessable press
hosts. The caps bound lookup in small domains (Section~\ref{lookup-exposure}).

Recency and formatting are shortcuts. Text from 2020--2022 differs in topic from
text from 2013--2016, model-only mirrors are newer still, and generated text
carries typography and layout that human excerpts lack. Uniform sampling over
unique documents puts about 44\% of the mass in 2020--2022, and each model-only
mirror is conditioned on its parent's topic and on what existed by the parent's
capture date. A versioned cohort generation profile then puts each generated
child into its parent's form: a parent-conditioned normaliser applies the
parent's quote, dash, range, and ellipsis conventions and paragraph structure, a
form screen rejects children whose text fails the human excerpts' prose screen,
and one window of 150 to 450 words, the human excerpts' own bound, applies to
the final documents of every class. Two gates are measured on a Common
Crawl-parented mock batch, each passing on its point estimate with its bootstrap
upper bound reported: an era-only adversary and a layout-and-typography
adversary must each gain at most 5\% of the reference's headroom.
Character-class shares, such as case, digits, non-ASCII characters, and emoji,
and an era probe's ability to tell parents from mirrors are reported as
diagnostics, not gated. Both gates pass on the mock batch, but the era-only
gate's upper bound reaches 5.04\% (Section~\ref{cc-stock-evidence}). It is
accepted on its point estimate for the testnet phase and re-measured before
mainnet live rewards on normalised, length-matched positives from the next
training round, with a cohort large enough to bring the upper bound below 5\%;
until then the bound is a stated limit. \(\sigma\), \(\delta\), \(w\), and the
headroom are re-derived on Common Crawl-parented batches. Common Crawl text is
evaluated privately and never redistributed. As project policy it may be sent to
hosted model APIs to generate mixed rewrites.

Testnet cohorts use a separate public stock: the unassigned remainder of the
evaluation-reserved ranges behind the sealed Test, about 240,000 dated public
documents ranging from Wikipedia and PLOS articles to forum and blog text, all
historical-human proxies (Section~\ref{human-evaluation-stock}). At 1,000
cases it lasts about three years ignoring domain balance, but seven of its
eighteen domains hold fewer than 350 documents, so domains bind first. When a
domain cannot fill its share, a later manifest reweights or retires it and
records the changed estimand; \(\sigma\), \(\delta\), and \(w\) are then
re-derived. When a stock cannot fill a batch, or the mainnet floor, no new
cohort opens and the last target stays live, paying the same clusters, until new
stock arrives. Every text in the public stock can be recognized
(Section~\ref{lookup-exposure}); this is an accepted pilot limit of testnet,
which reaches mainnet only through the reference testnet hands on
(Section~\ref{reference-model}).

Testnet keeps fresh stock for cohorts with outside entrants. A cohort whose
roster has no admitted outside entrant with a working artifact, such as a
rehearsal with only operator or synthetic entrants, draws its human documents
from already-exposed pools alone: 570 development human documents, 4,593 public
documents whose only earlier roles were the consumed Test, held-out, or
diagnostic sets, and any document an earlier testnet cohort drew. These may be
scored again as pure-human negatives without limit, never-drawn documents first
and none twice in one cohort, while their uses as rewrite sources keep the
limits above. Fresh stock is drawn only for a cohort with such an entrant. The
coordinator enforces the rule: it draws only from pools on its exposed list and
refuses any plan that names another. On-chain admission is not yet wired into
cohorts, so every testnet cohort so far has drawn exposed documents only. A
score on reused documents shows that an artifact runs and is scored, not how it
performs on unseen text.

\subsection{Lookup exposure and the canary}\label{lookup-exposure}

The human side of a batch is existing text, and a mixed document keeps the
unreplaced parts of its source byte for byte. An artifact that memorized human
text could look a passage up and compare it with the query, labelling
recognized text human and altered text AI, without detecting AI involvement.
Its gain is the covered documents' share of the reference's loss, roughly
coverage times headroom, and more where it targets domains the reference
handles worst. Admission rejects an explicit text index
(Section~\ref{a-artifact-execution}); an artifact later found to carry one has
invalid scores, and their signed correction must also address any promotion
they produced. Recognition stored in the parameters, including through
pretraining exposure, cannot be excluded. The protocol reduces lookup; it does
not prevent it.

On testnet the exposure is large. The whole public stock fits in a 4.1 MB
sentence Bloom filter, and its raw text in 327 MB, far below any artifact cap.
Recognizing one domain's documents alone, such as Wikipedia, would gain about
12\% of the reference's headroom on a 40/40/20 batch, more than the 1,000-case
margin. On mainnet the frame is web-scale, about 64 billion first-seen URLs, while an artifact at
the cap holds an estimated 40--72 MB of learned capacity: 160 million parameters
at the 2 to 3.6 bits per parameter that language models have been measured to
store~\cite{allenzhuknowledge,morrismemorize}. Against that capacity,
estimated gains are at most about 0.09\% of headroom under uniform sampling over
unique documents, 0.31--0.46\% if sampling followed captures, and 2.3--3.4\% if
every domain had an equal share. The share caps hold the total to about 0.5\%
against an attacker whose coverage of each stratum stays partial; one who
completes a capped stratum can reach about 1.1\%. These estimates can
understate coverage by up to 1.44 times, and the frame size is known only to
within a factor of two. Encoders pretrained on web text may already recognize
Common Crawl pages. A first red team found no membership signal in the
masked-language-model loss of two small encoders pretrained on data containing
Common Crawl, and a recognizer fine-tuned from \(R_0\)'s encoder recovered
8.5--14.7\% of the pages it was trained on at a 0.1\% false-positive rate, far
below the ideal index that the caps price; the test on ModernBERT-base's own
pretrained weights is still owed (Section~\ref{cc-stock-evidence}).

The defences are the artifact cap; breadth of sampling; keyed, secret passage
choice, which multiplies the keys an attacker must store by about 5.5, the
average number of valid passages per page; share caps for small or guessable
domains; text first published after a manifest freezes, such as new
parliamentary records or court opinions, as an option, which counts as unseen
once only; and a lookup canary.

The canary uses about 1,000 freshly written anchor documents a year, produced by
paid writers on a keystroke-logged editor, plus about \$1,000 a year of
US-government court filings bought from PACER that are not already in free
public archives. Anchors stay private, form one source group per writer, are
reused across audits, and are never scored in a cohort. They cost about \$8,900
a year, which keeps the human-data budget under \$10,000 a year. Every
candidate for promotion to reference (Section~\ref{reference-model}) is
audited: its mean paired utility difference against a frozen \(R_0\) on a fresh
20,000-document Common Crawl draw, minus the same difference on the anchors and
their rewrites, is flagged when it exceeds an honest-model offset by 3.09
standard errors. A first flag triggers an immediate re-audit in the same cohort
on a fresh, independent Common Crawl draw. A cohort's candidates share one
draw, so only a re-audit adds a second. At one audit a day, the draws take
about 7.3 million fresh documents a year, 25 times the human negatives of
4,000-case batches and about 1\% of the frame's WET text
(Section~\ref{cc-stock-evidence}), and scoring one model on a draw costs about
\$0.70. A single flag is only reported in the signed package. Two flags block the candidate's promotion, with no automatic
reward penalty: the next cluster head in rank order that survives the near-equal
drop (Section~\ref{reward-v1}) and passes the canary and
every other promotion condition is promoted instead, and credits are paid as
usual. A signed operator review may later clear a blocked candidate; clearing
does not change a past promotion. The canary has little power: with 1,000 anchors, a
close fine-tune of \(R_0\) needs a lookup gain of about 20--23\% of headroom to
be flagged with 80\% probability, and no achievable lookup by a very different
model is flagged. Honest models are falsely flagged in about 0.3--0.6\% of
audits. Prevention therefore does most of the work.

\subsection{AI detection quality}\label{detection}

Let \(a_i(t)\) be model \(i\)'s predicted AI-involvement probability for document
\(t\), its displayed AI confidence divided by 100 and the sum of its model-only
and mixed probabilities. The independent label \(y_t\) is one for either
model-only or mixed production and zero for documented human-only work under
the frozen labeling policy. Binary Brier utility is
\begin{equation}\label{eq:brier-utility}
 u_i(t)=1-(a_i(t)-y_t)^2.
\end{equation}
For \(G\) source groups, let \(\lambda_g\geq0\) sum to one and let within-group
weights \(\omega_{t\mid g}\geq0\) also sum to one. The self-weighting batch makes
these weights equal per document. The cohort score and its improvement over the
reference \(R_k\) in force are
\begin{equation}\label{eq:a-quality}
 S_i=\sum_{g=1}^{G}\lambda_g\sum_{t\in g}\omega_{t\mid g}u_i(t),
 \qquad
 q_i=\frac{S_i-S_{R_k}}{1-S_{R_k}},\quad S_{R_k}<1,
\end{equation}
where \(S_{R_k}\) is the reference's score on the same batch under the same
outcome rules. Improvement is divided by the reference's remaining headroom;
negative examples are not individually clipped away. Involvement probabilities
enter in parts per billion, so each utility is exact in units of \(10^{-18}\),
and a score is the floor of the batch mean on that grid. Cutoffs, bands, and the
near-equal test then compare these integers as exact rationals without further
rounding.

AI-involvement Brier is the only reward quantity. Share, attribution, span, and
the pure-human false-positive gates decide eligibility or promotion and carry
no reward weight. Evaluate share against measured fractional labels with source-balanced
absolute and squared error, coverage, and slices by rewrite size and procedure.
A mixed-history label does not imply 50\% AI share, and a fractional share target
cannot replace \(y_t\) in Equation~\ref{eq:brier-utility} while preserving its
involvement meaning. Span outputs are measured as described in
Section~\ref{span-outputs}.

Report three-class origin Brier, calibration, sensitivity at fixed human
false-positive rates, and population slices separately. Attribution reports
include retrieval accuracy, reciprocal rank, open-set rejection, and performance
on withheld writers and generators. Rejection means declining named attribution
as defined in Section~\ref{author-and-origin-detection}, including
\emph{Other human} and \emph{Other AI}. These comparisons use independent
reference and query works. A profile match cannot override the origin output,
and mixed work need not have one attributable author. Measure \emph{Other AI} on
documented generator roles and \emph{Other human} on documented human roles,
including both roles in mixed work. Neither fallback earns named-identity credit.
Both outputs are conditional, so a fallback for an absent role is not itself a
false attribution; mistaken involvement probabilities are scored through origin
measurements. Low predicted role probability never removes a documented
contributor from the attribution denominator. Gate thresholds, including both
false-positive tiers, are fixed on development evidence before the cohort.
Because the gates carry no reward weight, they are floors against unusable
outputs rather than improvement targets. That evidence includes \(R_0\)'s own
mixed-document share error of 15.5 points and attribution top-1 accuracy of
0.35 (Section~\ref{current-a-baseline}); the thresholds themselves are not yet
chosen (Section~\ref{development-and-launch-criteria}).

\subsection{The reference model}\label{reference-model}

The reward baseline is the current reference model \(R_k\). The sequence
starts at \(R_0\), the document model of the \textit{A} baseline release
candidate described in Section~\ref{existing-evidence-and-implementation-status},
published like every reference (Section~\ref{model-interfaces-and-starter-architectures}). Each cohort re-runs
\(R_k\), and its predecessor \(R_{k-1}\) when one exists, on that cohort's batch.
When a cohort's top-ranked eligible entrant, the head of its top cluster below,
improves on \(R_k\) by the margin \(\delta\), passes the strict false-positive
tier (Section~\ref{cohort-batch}), and is not blocked by the lookup canary
(Section~\ref{lookup-exposure}), that entrant becomes \(R_{k+1}\): the baseline
for the next cohort and the new reference, published with its cohort's
artifacts when the cohort ends. When it fails either check, the next retained cluster head in rank order that
passes every promotion condition is promoted instead, and credits are paid as
usual; a cluster removed by the near-equal drop is never a candidate.
Otherwise the reference does not change.
Promotion uses the same margin as eligibility, with no larger promotion
threshold. Because the margin uses the pooled standard error of seed retrains,
how often it promotes an equally good model depends on how entrants differ from
\(R_k\) (Table~\ref{tab:promotion-rates}). Entrants that differ from it as seed
retrains do share \(R_k\)'s scoring noise, which lowers the chance that one of
several clears the margin. A structurally different model has a much
larger paired standard error against \(R_k\): those of the other two measured
recipes and of V6 against \(R_0\) are about 13--19\% of headroom at 1,000 cases,
3.7 to 5.4 times \(\sigma\) (Section~\ref{cohort-mix-evidence}). Those models
are much weaker than \(R_0\), so they supply only this scale, and the table's
figures for structurally different entrants, taken at its midpoint of about
\(4.6\sigma\) under a normal approximation, are illustrative. Which kind of
roster the subnet will see is not yet measured. Once independently trained
models of equal quality compete, promotion rather than fallback is the expected
case, and the reference changes hands among equally good models. The same noise
can promote a worse model; the strict false-positive tier and the canary bound
human false positives and lookup, not quality.

\begin{table}[H]
\centering
\begin{tabularx}{\linewidth}{@{}llL@{}}
\toprule
Entrants & Quality against \(R_k\) & Chance that one clears the margin \\
\midrule
Ten, differing as seed retrains & equal & about 14\%; 21\% if their
differences were independent \\
One, structurally different & equal & about a third (0.29--0.36 across the
measured range) \\
Ten, structurally different & equal & about 80\% if they share \(R_k\)'s noise
as seed retrains do; 98\% if they share none \\
One, structurally different & half a margin below & about a quarter \\
One, structurally different & one margin below & about a fifth \\
One, structurally different & two margins below & about a tenth \\
\bottomrule
\end{tabularx}
\caption{Chance per cohort that the shared margin promotes an entrant no better
than \(R_k\).}
\label{tab:promotion-rates}
\end{table}

No rule bounds the resulting drift in the
reference's true quality, but four properties limit it. The promoted model is
the top-ranked entrant that passes every promotion condition, so a worse entrant
is promoted only when it outscores every better qualifying entrant on that
batch. Each cohort re-runs the reference on a fresh batch, so a winning draw
does not carry forward. A worse reference is easier to pass: against one a full
margin below the model it displaced, a structurally different entrant as good as
that model clears the margin about half the time rather than a third. And the
displaced model becomes \(R_{k-1}\), so a fallback cohort pays only entrants that
outscore it. False promotion and the drift it allows are stated limits of the
shared margin; their size on simulated rosters is measured before launch
(Section~\ref{protocol-checks-before-launch}).

Testnet opens a fresh record of cohorts and references, with \(R_0\) as its
first reference and no predecessor. One change there falls outside the promotion
rule: \(R_0\) may be replaced once by the next reference Slop Ninja trains,
retrained and recalibrated on normalised, length-matched positives
(Section~\ref{human-stock}), and only if that model passes every required
gate. The status page announces the replacement in advance. The
replacement takes \(R_0\)'s position, with no predecessor, and appears in the
next signed manifest, so the sequence of signed manifests is its record. Credit
is bound to its own cohort's package and target (Section~\ref{lagged-payout}),
so the replacement moves no earned reward. It sets the margin for later cohorts
and, while the launch burn lasts, the bar that ends it. A span package,
whether the release package or the first one that serves spans, has different
inference content from \(R_0\), so naming it as the reference would itself use
that one replacement, although its document outputs are identical; no span
package is named \(R_0\), and the replacement stays available. Mainnet opens
its own fresh record, whose first reference is the one testnet holds at that
time, with no predecessor and no separate one-time replacement. A testnet
promotion can rest partly on lookup, since testnet's stock is recognizable in
full, its strict tier is drawn from the same pools, and its canary has no
anchors. Mainnet scores that reference only on its own batches, where
recognizing testnet text gains nothing, so mainnet's first bar is the
reference's measured mainnet quality, not its testnet record. An inflated record
therefore leaves that bar lower than the record suggests: mainnet's launch burn
can end, and its first promotion occur, against a weaker model, which once
displaced anchors the fallback cutoffs until the next promotion. Before
mainnet, that reference is measured on Common Crawl and the anchors, and the
results are reported (Section~\ref{development-and-launch-criteria}).

Every entrant artifact is published after its cohort ends, and a copy of a
published artifact earns nothing. An entrant whose inference content equals that of any
published artifact is not executed. An executed entrant within the near-copy
distance of a published artifact (Section~\ref{reward-v1}) fails the public-copy
gate unless it scores more than half a band above that artifact. This applies to
\(R_k\) and \(R_{k-1}\), including \(R_0\) before its package is published, and to every
earlier entrant's artifact, the author's own included. A miner therefore earns again only with a model that differs
measurably from everything already published, such as an independent retrain or
a weight soup of published models, or that improves on its source by half the
promotion margin. The gate removes exact and recalibrated copies; it sets a
minimum price on a copy perturbed with deliberate noise, but cannot prevent one
or bound what it earns (Section~\ref{reward-v1}).

\subsection{From scores to credits}\label{reward-v1}

An entrant is \emph{gated} when its artifact qualification, its public-copy
gate, its share and attribution gates, and the per-entrant tier of its
pure-human false-positive gate pass, together with the span gate once a
manifest enforces it. Gate failures
earn no credit. The margin is
\[
 \delta=\max\bigl(0.01,\;2\sigma\bigr),
\]
where \(\sigma\) is the pooled paired standard error of \(q\) at the manifest's
batch size, frozen in the manifest. Under 20/55/25, \(\delta\) is about 0.070 at
1,000 cases and 0.034 at 4,000. Each manifest re-derives \(\sigma\) from current
measurements. One frozen \(\sigma\) is used rather than each entrant's own
paired standard error because it also sets the cohort's single band width, and
because it fixes every entrant's bar before any score is seen instead of making
the bar depend on how closely an entrant tracks \(R_k\)'s predictions.

Table~\ref{tab:reward-v1-parameters} lists the remaining parameters and their
translation from subnet 78. Subnet 78 states its band width and duplicate
tolerance in points of a normalized score. A literal 2.5-point band would span
about 40--48\% of the reference's headroom here, depending on the estimand, and
put almost every entrant in one band. \texttt{reward-v1} therefore keeps the
mechanism and the ratios between its parameters, and sets the band width from
the measured standard error.

\begin{table}[H]
\centering
\begin{tabularx}{\linewidth}{@{}lLL@{}}
\toprule
Parameter & Subnet 78 & \texttt{reward-v1} \\
\midrule
Band width \(w\) & 2.5 points of score & \(2\sigma(1-S_{R_k})\) in score units:
7\% of headroom at 1,000 cases (about 0.0044 at 0.062 headroom), 3.4\% at 4,000 \\
Near-equal tolerance \(\varepsilon\) & 1 basis point, \(w/250\) & \(w/250\), inclusive \\
Decay \(r\), across and within bands & 0.85 & 0.85 \\
Eligibility & \(S\geq0.99\,S_{\mathrm{inc}}\), the incumbent's score & gated and \(q\geq\delta\) \\
Duplicates & identical content credited to the first acceptance & copies of
published artifacts earn nothing; near-duplicate entrants share one credit
equally \\
Near-copy distance \(d\) & --- & 0.025 after isotonic recalibration, with a
\(w/2\) improvement exception \\
Native row & largest-remainder split of 65,535 & the same \\
\bottomrule
\end{tabularx}
\caption{Frozen \texttt{reward-v1} parameters, signed in each cohort manifest.}
\label{tab:reward-v1-parameters}
\end{table}

Two models are within the \emph{near-copy distance} when, after the best
rank-preserving recalibration, their involvement probabilities still differ by
a root-mean-square of at most \(d\), and the first does not score more than
\(w/2\) above the second, where \(w\) is the cohort's band width. The recalibration is fitted both ways by isotonic
regression over the documents on which both models returned a prediction, and
the smaller residual counts; two models with no common prediction are not near.
A temperature, bias, or any other change that preserves a model's ranking of
documents is therefore at distance zero, however far it moves the
probabilities. The manifest freezes \(d=0.025\). On 2,000 random 1,024-document
40/40/20 batches drawn from the consumed Test, two training seeds of one recipe
never came closer than 0.034, and seeds of the launch recipe never closer than
0.059.
A model built on a published one stays eligible when it beats it by more than
\(w/2\), which is half the promotion margin while \(2\sigma\geq0.01\).

The calculation proceeds in seven steps. Entrants are ordered by score, then by
finalized commitment block and extrinsic index, then by hotkey bytes; this
\emph{rank order} breaks every tie below.
\begin{enumerate}\tightlist
\item \emph{Cutoff.} If any gated entrant has \(q_i\geq\delta\), the cutoff is
\(T=S_{R_k}+\delta(1-S_{R_k})\). Otherwise, when \(R_{k-1}\) exists and some gated
entrant scores strictly above \(S_{R_{k-1}}\), the cohort is scored against the
fallback cutoff
\[
 C_N=S_{R_{k-1}}+0.15\,\bigl(\max_N-S_{R_{k-1}}\bigr),
\]
where \(\max_N\) is the top gated score in cohort \(N\). The fallback removes the
bottom 15\% of the score range that would earn against the previous reference,
so rewards keep flowing to the current cohort; the reference does not advance.
The anchor is \(R_{k-1}\) rather than \(R_k\) because no entrant has passed
\(R_k\) by the margin, so none is measurably different from it, and a cutoff at
\(S_{R_k}\) would admit or exclude entrants equal to it by batch noise. Anchoring
on the last reference the subnet passed by the margin lets independent models as
good as \(R_k\), including new models from \(R_k\)'s author, earn in a fallback
cohort, while copies of \(R_k\) fail the public-copy gate
(Section~\ref{reference-model}). For example, if \(R_{k-1}\) scores 0.9400 on
the batch and the top entrant 0.9500, the cutoff is 0.9415. An entrant just
short of the margin therefore earns in a fallback cohort but not when another
entrant passes it. In any other
case, including \(S_{R_k}=1\), where the band width is zero even when \(R_{k-1}\)
exists, the cohort produces no new target: the previous target stays live, and
if there has never been one the row sends the whole pool to the owner hotkey's
UID, which burns it.
\item \emph{Clusters.} Walk the gated entrants at or above the cutoff in rank
order. The first entrant not yet assigned heads a cluster that also contains
every later unassigned entrant with identical inference content or within the
near-copy distance of the head on the batch. The cluster takes its head's score
and position; there is no chaining through other members.
\item \emph{Near-equal drop.} Walk the cluster heads in rank order and drop a
cluster whose score is at most \(\varepsilon\) below the last retained head's
score; of two independent clusters with exactly equal scores, the later in rank
order is therefore dropped. Comparison is always with the last retained head,
never with a dropped one, and happens before banding. A dropped cluster earns
nothing.
\item \emph{Bands.} Place each retained cluster in band
\(b=\min(\lfloor S/w\rfloor,\lceil 1/w\rceil-1)\) on the absolute score grid.
Rank occupied bands \(\kappa=0,\dots,B-1\) from the top, and rank the clusters
within band \(\kappa\) as \(\mu=0,\dots,n_\kappa-1\) in rank order.
\item \emph{Credit.} A cluster earns
\[
 c=\frac{r^{\kappa}}{\sum_{j<B}r^{j}}\cdot\frac{r^{\mu}}{\sum_{j<n_\kappa}r^{j}},
\]
so the credits sum to one.
\item \emph{Equal split.} Each of a cluster's \(m\) members receives \(c/m\).
\item \emph{Row.} Apportion 65,535 units across members by largest remainder:
each takes the integer part of its share, and leftover units go to the largest
fractional remainders, with ties in rank order. At each tempo, map each member's
hotkey and registration generation to its current UID in the validator's
finalized registration snapshot. Units of a hotkey that is no longer registered
under that generation go to the UID of the subnet owner hotkey in the same
snapshot, which burns them (Section~\ref{unallocated-rewards}). The chain stores
a weight row scaled so that its largest entry is 65,535, so validators commit,
and compare each reveal with, the row in that form, computed exactly with ties
rounded half to even; the canonical row and its digest remain the audited
record.
\end{enumerate}

Because an entrant short of the margin earns only in a fallback cohort, how
often an equally good entrant earns depends on the roster. With ten entrants
equal to \(R_k\) that resemble its seed retrains, 79--86\% of cohorts fall back,
and in those an entrant earns whenever it reaches \(C_N\). Among ten
structurally different equal entrants, each earns only when it clears the margin
itself, about a third of the time, or in the 2--20\% of cohorts that fall back:
in a third to at most a half of cohorts, chosen by batch noise.
A copy of a published artifact earns nothing, so each artifact takes that draw
once, and equally good entrants earn equally only in expectation over many
artifacts.

With \(r=0.85\), a lone cluster receives the whole pool, and two clusters receive
54.1\% and 45.9\% whether they share a band or not, which the row encodes as
35,424 and 30,111 of its 65,535 units. A crowded band divides its
share: with two clusters in the top band and one below, they receive 29.2\%,
24.8\%, and 45.9\%. With \(r=17/20\), each credit is already a fraction in lowest
terms, and the package publishes every retained cluster's exact credit in that
form. All arithmetic uses exact integers and rationals, so every validator
derives byte-identical rows. A roster holds at most 256 entrants, the subnet's
UID limit, and because \(\varepsilon=w/250\), at most 250 retained clusters fit
in one band.

A band's share depends on its rank among occupied bands, not on how many
clusters it holds, so credit is not monotone in score across a band boundary. A
cluster alone in the lower of two occupied bands receives 45.9\% however many
clusters share the band above, 2.2 times the best of three there and 3.1 times
the best of five; alone in the third of three bands, it receives 28.1\%. Subnet
78's 2.5-point bands seldom divide a roster, so its experience does not cover
this case. A band of two standard errors often does: the paired difference
between two equally good entrants has a standard error of half a band, and the
grid is fixed in absolute score, so where a roster divides depends on the batch. A score
difference below the noise can therefore reverse the order of pay. The split is
kept from the ported mechanism, and what it buys is the same property: a band's
share does not shrink as other bands fill, so a crowd of entries in one band,
such as variants just beyond the near-copy distance, cannot dilute or push down
the pay of entrants in other bands, as it would under a geometric split by rank
alone. Its non-monotonicity is a stated limit; the strategy it opens is
described below.

The near-equal drop is kept from subnet 78, where it removes near-equal
duplicate entries. Here the content check, the public-copy gate, and
clustering remove copies before it runs. Its remaining benefits are fidelity to
the ported mechanism and a partial backstop against an entrant that escapes the
near-copy distance yet scores within \(\varepsilon\) of a retained head: the
drop cannot tell a copy from an original and removes whichever ranks lower. Its
main effect falls on independent entrants whose scores tie a retained head by
chance. Because
\(\varepsilon\) is \(1/125\) of the paired standard error at every batch size,
two equally good independent retrains fall within it in about 0.6\% of cohorts,
and among ten such entrants above the cutoff some cluster is dropped in about a
quarter of cohorts. That cost is a stated limit, like the false-promotion rate.

The near-copy distance removes exact and recalibrated copies and sets a minimum
price on a perturbed copy; it neither prevents one nor bounds what one earns. To
clear distance \(d\), a copy must add
variation that no rank-preserving recalibration explains. Variation of that size
that carries no label information costs about \(d^2\approx6\times10^{-4}\) of
expected score, less any gain also available by recalibrating the reference,
about \(3.7\times10^{-4}\) for \(R_0\), so the guaranteed floor is about
\(2.6\times10^{-4}\). The cheapest perturbed copy of \(R_0\), standing in for any
published reference, found on the consumed Test adds noise to \(R_0\)'s ranking
of documents after its best recalibration.
Pooling two independent measurements of 128 noise realizations each, on
40/40/20 batches, it cleared the gate on the batch and against \(R_0\)'s probe
fingerprint together in about 89\% to 99\% of 1,000-case cohorts, for an
average of about \(3.5\times10^{-4}\) to \(4.8\times10^{-4}\) of expected score:
0.6\% to 0.8\% of the reference's headroom, or 0.07 to 0.10 of a band at that
measurement's \(\sigma=0.040\), about a tenth of a band. At 4,000
cases the batch comparison alone is cleared nearer the floor, but the
1,024-document probe comparison with \(R_k\)'s fingerprint
(Section~\ref{public-copies}) then binds: clearing both in about 99\% of cohorts
costs about \(4.7\times10^{-4}\), about a fifth of a band. That cheapest copy
stays in the reference's band, takes a within-band slot in a fallback cohort, and
must be rebuilt each cohort because its own artifact is published. Once a
predecessor exists, fallback is the expected case while entrants
resemble seed retrains of the reference: ten such entrants equal to \(R_k\)
promote none about 86\% of the time, although structurally different equal
entrants usually promote one (Section~\ref{reference-model}). A tolerance
large enough to price an escape at a full band would be
about 0.07, near the distance between independent retrains of the launch
recipe, so no prediction tolerance can price it. With \(r=0.85\), an honest
cluster that shares a band with one such copy receives 54.1\% of the band's
share. Clustering uses the same distance, so a sybil variant of an entrant
costs the same to separate from its cluster.

Adding more noise escapes just as surely and lowers a copy's score by any chosen amount,
and because pay is not monotone across bands, a lower score can pay more. A
fallback cohort admits scores down to \(C_N\), 85\% of the way from the top
score to \(S_{R_{k-1}}\). When \(R_k\) genuinely improved on \(R_{k-1}\) by the
margin, about a band in score units, that range reaches most of a band below the
entrants near \(S_{R_k}\). A copy of \(R_k\) degraded into it, with noise well
beyond \(d\), can land alone in a band below theirs and take 45.9\% of the pool
if they occupy one band, or 28.1\% if they occupy two. In a promotion cohort, an
entrant whose model clears the margin by more than a band can do the same with a
weaker variant under a second hotkey, and any entrant can submit less than its
best model. Such an entry aims blind: the batch, the grid's position, and the
cutoff are unknown at commitment, and a miss earns nothing below the cutoff or a
low slot in a crowded band. Its expected credit against realistic rosters has
not been measured, and that measurement precedes launch
(Section~\ref{protocol-checks-before-launch}). Where it lands alone, it can earn
more than each honest cluster above it; this is a stated limit of the ported
split.

Weight soups of published artifacts are new models. A soup that is not within
the near-copy distance of a published artifact earns normally, and the protocol
does not look for one: averaging the weights of fine-tunes of one base, like
further training of a published model, is legitimate work. A submission must
still be one allow-listed graph, so an average of two separate encoders' outputs
cannot be submitted. Such an output average is the only blend measured: on the
consumed Test, a 50/50 average of two published seeds' outputs scored about 0.2
to 0.3 of a band above \(R_0\) at 1,000 cases and cleared the gate against both
sources, on the batch and on the probe, in every batch measured. A soup earns in
a fallback cohort like any model at its score, and one that beats \(R_k\) by the
margin can be promoted.

\subsection{Published artifacts and the fingerprint probe}\label{public-copies}

Exact copies are found by inference-content digest against every artifact ever
published, a check every validator repeats from the published packages: an
entrant fails the content check exactly when its digest equals that of a
published artifact, and such an entrant is never executed. Re-running every
earlier artifact on each batch would grow without bound, so near copies are also
measured on a fixed private probe of 1,024 documents, drawn and generated like a
batch. Near copies of \(R_k\) and \(R_{k-1}\) are measured on the cohort batch,
where both references already run, and on the probe against their
fingerprints. Near copies of every earlier artifact are measured on the probe.
A near copy of \(R_k\) is recorded before one of \(R_{k-1}\), and either before a
probe match.

Every executed entrant also runs on the probe once, whatever its score and other
gates. No score in an entrant's own cohort bounds what a copy of it can earn:
\(R_{k-1}\) can outscore \(R_k\) on a batch, and a copy meets new documents in a
later cohort. The entrant's probe predictions are compared with the recorded
probe predictions of every earlier published artifact, \(R_k\) and \(R_{k-1}\)
included, and become its own fingerprint once its artifact is published. Each
package publishes every entrant's exact comparisons with both references'
fingerprints, and validators check that a near one fails the gate, so a copy of
a reference must escape both the batch and the probe. When entrants with
identical inference content return different probe predictions, the
earliest-committed entrant's become the fingerprint. Each artifact therefore
runs on the probe only in its own cohort. An artifact that was never executed
carries no fingerprint and is protected by its content digest alone.

The probe carries no reward weight. Its labels serve only the \(w/2\) exception
against fingerprinted artifacts. Its human texts are withdrawn from the hidden
stock and never scored. The probe and the fingerprints stay private, and each
manifest binds at least one probe version, of at most 65,536 documents, by
version and digest. A model that could recognize probe documents could look
different there alone, so the probe follows the batch design. A new version is
drawn when that design changes materially or the probe may have been disclosed.
A replaced version stays in use, and entrants run on it too, while it holds the
fingerprint of any artifact not fingerprinted again on a later version. Both
references are fingerprinted on every version in use. A disclosed version
leaves use at once, and earlier artifacts recorded only on it keep exact-digest
protection only.

The probe adds 1,024 inferences per executed entrant and probe version in use.
That roughly doubles an entrant's scoring at 1,000 cases and adds a quarter at
4,000. Each package also reports, as exact integers, the smallest distance
between two cluster heads, between \(R_k\) and any gated entrant on the batch, and
between any gated entrant and any fingerprinted artifact on the probe, so a
later manifest can lower \(d\) before independent models approach it.

\subsection{Failures and the result package}\label{aggregation-and-failure-policy}

An attributable failed or missing model result has zero utility and retains its
assignment weight. A Container crash or out-of-memory failure is an
infrastructure error until it reproduces on a fresh instance, in
single-document isolation, under the frozen resource limits; a crash that
reproduces counts against the artifact, and those documents score zero utility
for that entrant only. A task whose every attempt ended in an infrastructure
error blocks the package: the cohort pauses for reconciliation under the same
manifest, and the previous target stays live. Operator downtime never becomes a
miner failure. If either reference model fails to return a prediction, the
package is held for a signed correction rather than scored against a damaged
baseline. Documents outside the input profile never enter the batch, so no
entrant can gain from a smaller successful subset. Later \textit{B}
training-service obligations are absent from an \textit{A}-only launch; no
fictitious requests or \textit{B} availability score is needed to make an
\textit{A} entrant eligible.

The signed result package binds the cohort and revision, the commitments with
their finalized block and extrinsic index, the artifact digests, credited
hotkeys and registration generations, coverage and outcome accounting, each
failed gate with its reason code, exact per-miner scores, the reference scores
\(S_{R_k}\) and \(S_{R_{k-1}}\), the cluster assignments, each scored entrant's
public-copy result, with its exact near-copy statistics against \(R_k\) and
\(R_{k-1}\) on the batch and against their fingerprints on the probe, and the
matched published artifact of any failure, the statistics behind each
distance-based cluster membership, the published-artifact index root and count,
the distance monitor, every \texttt{reward-v1} parameter, the cutoff mode, each
retained cluster's exact credit, and every member's integer units. A gate with
no recorded failure passed, and no gate measurement is published. An entrant
that fails artifact qualification is never executed and carries no score. After
the cohort closes the coordinator also publishes the artifacts and their
commitment openings. The package binds the R2 receipt index's root and count
before validators commit. Validators verify the full retained index, counts,
and receipt digests as specified in Section~\ref{cohorts}. Compact numerical
inputs and calculated rows survive the detailed receipts' retention window.

The frozen calculation specifies exact arithmetic, operation order,
tie-breaking, rounding, normalization, and integer chain encoding. It binds the
registration identities used to map earned credit to UIDs; no chain runtime
version is an input. Identical inputs must produce identical canonical rows on
Linux and macOS; implementation-dependent floating-point drift is not an
accepted tolerance. A registration change that requires a different conversion
needs an authenticated compatible delivery revision; a chain runtime upgrade
never does. Neither can silently change the scores or let a validator improvise
its own normalization.

Published score components are aggregate numerical inputs, never individual
document predictions or hidden labels. Cluster assignments and near-copy
statistics are coordinator measurements on hidden predictions, which validators
trust like the scores. Validators check that each recorded gate and cluster
decision follows from them and from the published scores, that every entrant's
probe comparison includes both references, that an entrant fails the content
check exactly when its content digest equals that of a published artifact, and
that each closed cohort adds one fingerprint for every executed entrant. They
recompute each credit and every member's units and refuse any difference.
Retain private evidence under access controls sufficient for authorized audit;
public record verification does not prove that the coordinator executed the
evaluation honestly. Corrections retain the original record and account for
already-submitted commitments and weights. The public inputs also make the row
derivable before reveal. Native commit--reveal therefore supplies the chain's
submission protocol without preventing off-chain copying or proving that each
validator did the calculation.

\subsection{Mechanism mapping and unallocated rewards}\label{unallocated-rewards}

The documented chain supports at most two native mechanisms with a shared UID
space, and under Yuma Consensus the UID limit times the mechanism count may not
exceed 256~\cite{mechanisms}. Launch uses one for \textit{A}. If \textit{C} later
activates, a candidate mapping uses native mechanism 0 for \textit{A} quality
and any paid-use credit, and mechanism 1 for grammar contributions, which
halves the subnet's UID slots to 128. \textit{B} needs a revised mapping before
activation; three logical workstreams do not create a third native slot.
\textit{C} may activate before \textit{B}; \textit{B}'s revised mapping must
then also place \textit{C}'s contribution credit.

Let \(c\) be the fraction of total miner emissions assigned to \textit{C} and
\(f\) the total fraction assigned to paid \textit{A} use, with
\(0\leq c<1\), \(f\geq0\), and \(c+f\leq1\). For separately normalized
quality, usage, and new grammar-credit distributions \(Q_i,I_i,K_i\), the
intended allocations before reward caps and the handling of zero-work usage
pools are
\begin{align}\label{eq:mechanism-allocation}
 W^A_i&=(1-c-f)Q_i+fI_i,\\
 W^C_i&=cK_i.
\end{align}
Here \(Q_i\) is the \texttt{reward-v1} credit. The native \textit{A} mechanism has
share \(1-c\); its submitted distribution divides \(W^A_i\) by \(1-c\). Launch
sets \(c=0\) and \(f=0\). A nonzero \(f\) requires the attribution and abuse
controls in Section~\ref{paid-asynchronous-jobs-and-on-chain-coordination}.

Unallocated credit goes to the UID of the subnet owner hotkey, which the chain
burns. Every validator reads the owner hotkey and its UID from the finalized
snapshot it uses for registrations, at every tempo, and records that hotkey in
the row. No manifest, package, or target names the owner hotkey, so an
owner-hotkey rotation moves the burn destination at the next tempo without a
new target or authorization. If units must burn while the finalized owner
hotkey holds no UID, that tempo's row is deferred rather than sent to any other
UID; a row with nothing to burn is still derived.

The burn has a documented native cost: before normalizing across subnets and
applying its emission gate, the chain multiplies each subnet's TAO-side emission
share by one minus the fraction of the last tempo's miner incentive that went to
owner hotkeys, recycled or burned~\cite{btemissions}. A partial burn, such as a
departed hotkey's units, shrinks the share accordingly. While the row sends the
whole pool to the owner hotkey's UID, as at launch, the subnet receives no
TAO-side emission and hence no TAO or alpha pool injection. Under the same
documented rules its participant-side alpha emission continues to accrue and is
divided by the emission weight vector described below: the miner slot follows
the row to the owner hotkey's UID and burns, while the validator and owner
slots are paid as configured. The burn regime therefore removes miner pay and
the TAO-side share, not validator or owner alpha. The native-chain rehearsal
(Section~\ref{protocol-checks-before-launch}) must confirm these effects.
Off-chain allocations cannot guarantee precise realized payouts, freeze native
emissions, or claw back rewards already distributed.

The split of subnet emissions among owner, validators, and miners is the
subnet's emission weight vector. The subnet owner may change it at most once a
week, and root at any time, through the chain's subnet-emission-weights call,
within chain bounds: the owner slot at most 33\% (and at least 18\% for a
subnet created through the chain's lease mechanism), the miner slot at least 33\%, and the
validator slot at most 50\% and above root stakers' baseline
share~\cite{btemissionweights}. The mainnet split is not selected. The
testnet subnet (netuid 601) uses an operator configuration that places the
validator share just above that baseline, its chain minimum, to exercise the
owner's path; it implies nothing about mainnet.
